\documentclass[12pt]{ximera}
\typeout{Start loading xmPreamble.tex}%

% Add here extra macro's that are loaded automatically by all documents of claas 'ximera' or 'xourse' in this repo

%%
%%  Example:
%%
% \newcommand{\R}{\mathbb{R}}
\usepackage{tikz}
\usetikzlibrary{positioning, arrows.meta}
\usepackage{multicol}
\usepackage{enumitem}
\usepackage{caption}
\usepackage{amsmath}
\usepackage{amsthm}
\usepackage{fontenc}

%% ---------------------------------------------------------------
%% Shared worksheet macros.
%%
%% These came from the Summer 2026 worksheet set, where each worksheet carried its
%% own copy. They have to live here instead: a xourse inlines every activity into a
%% single document, so duplicate \newcommand definitions across activities collide,
%% and the xourse gobbles \input so a per-topic macro file never loads.
%%
%% They use \providecommand, NOT \newcommand. ximera.cls loads this file automatically,
%% and every activity also carries an explicit \typeout{Start loading xmPreamble.tex}%

% Add here extra macro's that are loaded automatically by all documents of claas 'ximera' or 'xourse' in this repo

%%
%%  Example:
%%
% \newcommand{\R}{\mathbb{R}}
\usepackage{tikz}
\usetikzlibrary{positioning, arrows.meta}
\usepackage{multicol}
\usepackage{enumitem}
\usepackage{caption}
\usepackage{amsmath}
\usepackage{amsthm}
\usepackage{fontenc}

%% ---------------------------------------------------------------
%% Shared worksheet macros.
%%
%% These came from the Summer 2026 worksheet set, where each worksheet carried its
%% own copy. They have to live here instead: a xourse inlines every activity into a
%% single document, so duplicate \newcommand definitions across activities collide,
%% and the xourse gobbles \input so a per-topic macro file never loads.
%%
%% They use \providecommand, NOT \newcommand. ximera.cls loads this file automatically,
%% and every activity also carries an explicit \typeout{Start loading xmPreamble.tex}%

% Add here extra macro's that are loaded automatically by all documents of claas 'ximera' or 'xourse' in this repo

%%
%%  Example:
%%
% \newcommand{\R}{\mathbb{R}}
\usepackage{tikz}
\usetikzlibrary{positioning, arrows.meta}
\usepackage{multicol}
\usepackage{enumitem}
\usepackage{caption}
\usepackage{amsmath}
\usepackage{amsthm}
\usepackage{fontenc}

%% ---------------------------------------------------------------
%% Shared worksheet macros.
%%
%% These came from the Summer 2026 worksheet set, where each worksheet carried its
%% own copy. They have to live here instead: a xourse inlines every activity into a
%% single document, so duplicate \newcommand definitions across activities collide,
%% and the xourse gobbles \input so a per-topic macro file never loads.
%%
%% They use \providecommand, NOT \newcommand. ximera.cls loads this file automatically,
%% and every activity also carries an explicit \input{../xmPreamble}, so this file is
%% read twice. That was harmless while it held only \usepackage lines (LaTeX skips a
%% package it has already loaded) but a second \newcommand is a hard error.
%%
%%   \blank[width]        fill-in-the-blank rule (default 2.5cm)
%%   \showwork            "show and explain your work" reminder
%%   \showworkline        ... plus which schedule line was used
%%   \showworkyear        ... plus which year's schedule was used
%%   \schedhead{...}      centred bold caption above a tiered-rate table
%%   \samesquares[scale]{left}{right}   two equal-area squares, labelled
%%   \samecubes[scale]{left}{right}     two equal-volume cubes, labelled
%% ---------------------------------------------------------------

\providecommand{\blank}[1][2.5cm]{\underline{\hspace{#1}}}
\providecommand{\showwork}{\textbf{REMEMBER TO SHOW AND EXPLAIN YOUR WORK.}}
\providecommand{\showworkline}{\textbf{REMEMBER TO SHOW AND EXPLAIN YOUR WORK.
  Explanation should include how and why you know which line of the
  schedule was used to calculate this amount.}}
\providecommand{\showworkyear}{\begin{sloppypar}\textbf{REMEMBER TO SHOW AND
  EXPLAIN YOUR WORK. Explanation should include how and why you know which
  line of the year's tax schedule was used to calculate this tax
  amount.}\end{sloppypar}}
\providecommand{\schedhead}[1]{\begin{center}\textbf{#1}\end{center}}

\providecommand{\samesquares}[3][1]{%
\begin{center}
{\footnotesize These squares are the \textbf{same size}.}

\vspace{1.5mm}
\begin{tikzpicture}[scale=#1,every node/.style={scale=#1}]
  \draw (0,0) rectangle (2.4,2.4);
  \node[anchor=north west] at (0.15,2.25) {Area:};
  \node[anchor=east]  at (-0.15,1.2) {#2};
  \node[anchor=north] at (1.2,-0.15) {#2};
  \begin{scope}[xshift=4.6cm]
    \draw (0,0) rectangle (2.4,2.4);
    \node[anchor=north west] at (0.15,2.25) {Area:};
    \node[anchor=east]  at (-0.15,1.2) {#3};
    \node[anchor=north] at (1.2,-0.15) {#3};
  \end{scope}
\end{tikzpicture}
\end{center}}

\providecommand{\samecubes}[3][1]{%
\begin{center}
{\footnotesize These cubes are the \textbf{same size}.}

\vspace{1.5mm}
\begin{tikzpicture}[scale=#1,every node/.style={scale=#1}]
  \foreach \dx in {0,5.2} {
    \begin{scope}[xshift=\dx cm]
      \draw (0,0) rectangle (2.0,2.0);
      \draw (0,2.0) -- (0.65,2.6) -- (2.65,2.6) -- (2.0,2.0);
      \fill[black!12] (2.0,0) -- (2.65,0.6) -- (2.65,2.6) -- (2.0,2.0) -- cycle;
      \draw (2.0,0) -- (2.65,0.6) -- (2.65,2.6) -- (2.0,2.0);
      \node[anchor=north west] at (0.1,1.9) {Volume:};
    \end{scope}
  }
  \node[anchor=east]  at (-0.15,1.0) {#2};
  \node[anchor=north] at (1.0,-0.15) {#2};
  \node[anchor=west]  at (2.25,0.4) {#2};
  \node[anchor=east]  at (5.05,1.0) {#3};
  \node[anchor=north] at (6.2,-0.15) {#3};
  \node[anchor=west]  at (7.45,0.4) {#3};
\end{tikzpicture}
\end{center}}

%% NOTE: do NOT add \usepackage to individual activity .tex files.
%% When a xourse inlines an activity it gobbles \documentclass and \input, but a bare
%% \usepackage still executes -- after \begin{document} -- and errors out.
%% All shared packages and macros belong here.
%%
%% ximera.cls already loads: amsmath, amssymb, amsthm, cancel, enumitem, graphicx,
%% hyperref, listings, pgfplots, tikz, url, xcolor. No need to re-load those.
, so this file is
%% read twice. That was harmless while it held only \usepackage lines (LaTeX skips a
%% package it has already loaded) but a second \newcommand is a hard error.
%%
%%   \blank[width]        fill-in-the-blank rule (default 2.5cm)
%%   \showwork            "show and explain your work" reminder
%%   \showworkline        ... plus which schedule line was used
%%   \showworkyear        ... plus which year's schedule was used
%%   \schedhead{...}      centred bold caption above a tiered-rate table
%%   \samesquares[scale]{left}{right}   two equal-area squares, labelled
%%   \samecubes[scale]{left}{right}     two equal-volume cubes, labelled
%% ---------------------------------------------------------------

\providecommand{\blank}[1][2.5cm]{\underline{\hspace{#1}}}
\providecommand{\showwork}{\textbf{REMEMBER TO SHOW AND EXPLAIN YOUR WORK.}}
\providecommand{\showworkline}{\textbf{REMEMBER TO SHOW AND EXPLAIN YOUR WORK.
  Explanation should include how and why you know which line of the
  schedule was used to calculate this amount.}}
\providecommand{\showworkyear}{\begin{sloppypar}\textbf{REMEMBER TO SHOW AND
  EXPLAIN YOUR WORK. Explanation should include how and why you know which
  line of the year's tax schedule was used to calculate this tax
  amount.}\end{sloppypar}}
\providecommand{\schedhead}[1]{\begin{center}\textbf{#1}\end{center}}

\providecommand{\samesquares}[3][1]{%
\begin{center}
{\footnotesize These squares are the \textbf{same size}.}

\vspace{1.5mm}
\begin{tikzpicture}[scale=#1,every node/.style={scale=#1}]
  \draw (0,0) rectangle (2.4,2.4);
  \node[anchor=north west] at (0.15,2.25) {Area:};
  \node[anchor=east]  at (-0.15,1.2) {#2};
  \node[anchor=north] at (1.2,-0.15) {#2};
  \begin{scope}[xshift=4.6cm]
    \draw (0,0) rectangle (2.4,2.4);
    \node[anchor=north west] at (0.15,2.25) {Area:};
    \node[anchor=east]  at (-0.15,1.2) {#3};
    \node[anchor=north] at (1.2,-0.15) {#3};
  \end{scope}
\end{tikzpicture}
\end{center}}

\providecommand{\samecubes}[3][1]{%
\begin{center}
{\footnotesize These cubes are the \textbf{same size}.}

\vspace{1.5mm}
\begin{tikzpicture}[scale=#1,every node/.style={scale=#1}]
  \foreach \dx in {0,5.2} {
    \begin{scope}[xshift=\dx cm]
      \draw (0,0) rectangle (2.0,2.0);
      \draw (0,2.0) -- (0.65,2.6) -- (2.65,2.6) -- (2.0,2.0);
      \fill[black!12] (2.0,0) -- (2.65,0.6) -- (2.65,2.6) -- (2.0,2.0) -- cycle;
      \draw (2.0,0) -- (2.65,0.6) -- (2.65,2.6) -- (2.0,2.0);
      \node[anchor=north west] at (0.1,1.9) {Volume:};
    \end{scope}
  }
  \node[anchor=east]  at (-0.15,1.0) {#2};
  \node[anchor=north] at (1.0,-0.15) {#2};
  \node[anchor=west]  at (2.25,0.4) {#2};
  \node[anchor=east]  at (5.05,1.0) {#3};
  \node[anchor=north] at (6.2,-0.15) {#3};
  \node[anchor=west]  at (7.45,0.4) {#3};
\end{tikzpicture}
\end{center}}

%% NOTE: do NOT add \usepackage to individual activity .tex files.
%% When a xourse inlines an activity it gobbles \documentclass and \input, but a bare
%% \usepackage still executes -- after \begin{document} -- and errors out.
%% All shared packages and macros belong here.
%%
%% ximera.cls already loads: amsmath, amssymb, amsthm, cancel, enumitem, graphicx,
%% hyperref, listings, pgfplots, tikz, url, xcolor. No need to re-load those.
, so this file is
%% read twice. That was harmless while it held only \usepackage lines (LaTeX skips a
%% package it has already loaded) but a second \newcommand is a hard error.
%%
%%   \blank[width]        fill-in-the-blank rule (default 2.5cm)
%%   \showwork            "show and explain your work" reminder
%%   \showworkline        ... plus which schedule line was used
%%   \showworkyear        ... plus which year's schedule was used
%%   \schedhead{...}      centred bold caption above a tiered-rate table
%%   \samesquares[scale]{left}{right}   two equal-area squares, labelled
%%   \samecubes[scale]{left}{right}     two equal-volume cubes, labelled
%% ---------------------------------------------------------------

\providecommand{\blank}[1][2.5cm]{\underline{\hspace{#1}}}
\providecommand{\showwork}{\textbf{REMEMBER TO SHOW AND EXPLAIN YOUR WORK.}}
\providecommand{\showworkline}{\textbf{REMEMBER TO SHOW AND EXPLAIN YOUR WORK.
  Explanation should include how and why you know which line of the
  schedule was used to calculate this amount.}}
\providecommand{\showworkyear}{\begin{sloppypar}\textbf{REMEMBER TO SHOW AND
  EXPLAIN YOUR WORK. Explanation should include how and why you know which
  line of the year's tax schedule was used to calculate this tax
  amount.}\end{sloppypar}}
\providecommand{\schedhead}[1]{\begin{center}\textbf{#1}\end{center}}

\providecommand{\samesquares}[3][1]{%
\begin{center}
{\footnotesize These squares are the \textbf{same size}.}

\vspace{1.5mm}
\begin{tikzpicture}[scale=#1,every node/.style={scale=#1}]
  \draw (0,0) rectangle (2.4,2.4);
  \node[anchor=north west] at (0.15,2.25) {Area:};
  \node[anchor=east]  at (-0.15,1.2) {#2};
  \node[anchor=north] at (1.2,-0.15) {#2};
  \begin{scope}[xshift=4.6cm]
    \draw (0,0) rectangle (2.4,2.4);
    \node[anchor=north west] at (0.15,2.25) {Area:};
    \node[anchor=east]  at (-0.15,1.2) {#3};
    \node[anchor=north] at (1.2,-0.15) {#3};
  \end{scope}
\end{tikzpicture}
\end{center}}

\providecommand{\samecubes}[3][1]{%
\begin{center}
{\footnotesize These cubes are the \textbf{same size}.}

\vspace{1.5mm}
\begin{tikzpicture}[scale=#1,every node/.style={scale=#1}]
  \foreach \dx in {0,5.2} {
    \begin{scope}[xshift=\dx cm]
      \draw (0,0) rectangle (2.0,2.0);
      \draw (0,2.0) -- (0.65,2.6) -- (2.65,2.6) -- (2.0,2.0);
      \fill[black!12] (2.0,0) -- (2.65,0.6) -- (2.65,2.6) -- (2.0,2.0) -- cycle;
      \draw (2.0,0) -- (2.65,0.6) -- (2.65,2.6) -- (2.0,2.0);
      \node[anchor=north west] at (0.1,1.9) {Volume:};
    \end{scope}
  }
  \node[anchor=east]  at (-0.15,1.0) {#2};
  \node[anchor=north] at (1.0,-0.15) {#2};
  \node[anchor=west]  at (2.25,0.4) {#2};
  \node[anchor=east]  at (5.05,1.0) {#3};
  \node[anchor=north] at (6.2,-0.15) {#3};
  \node[anchor=west]  at (7.45,0.4) {#3};
\end{tikzpicture}
\end{center}}

%% NOTE: do NOT add \usepackage to individual activity .tex files.
%% When a xourse inlines an activity it gobbles \documentclass and \input, but a bare
%% \usepackage still executes -- after \begin{document} -- and errors out.
%% All shared packages and macros belong here.
%%
%% ximera.cls already loads: amsmath, amssymb, amsthm, cancel, enumitem, graphicx,
%% hyperref, listings, pgfplots, tikz, url, xcolor. No need to re-load those.

\title{Exponential Decay: Predicting the Model Before You Roll}
\begin{document}
\begin{abstract}
Extending the dice-decay experiment to 4-, 8-, and 12-sided dice, and then the payoff:
you can write down a decent model \emph{without rolling anything at all}, because the
decay rate is just the probability of surviving a roll. Comparing the theoretical models
against the ones the spreadsheet fitted to real data.
\end{abstract}
\maketitle

\section*{The Class Models}

Each group ran the same procedure --- roll a pile of dice, remove every die showing a
``1,'' repeat until none remain --- but with different numbers of sides. Pooling the
class data, the spreadsheet produced these models, along with the group's actual starting
count:

\begin{center}
\begin{tabular}{lccc}
\hline
Dice & Model, $y = a \cdot e^{bx}$ & Actual starting \# & $R = e^{b}$ \\
\hline
6 sided  & $y = 136 \cdot e^{-0.181x}$   & 129 & 0.8344 \\
4 sided  & $y = 44.1 \cdot e^{-0.278x}$  & 44  & 0.7573 \\
8 sided  & $y = 32.8 \cdot e^{-0.131x}$  & 35  & 0.8772 \\
12 sided & $y = 33.4 \cdot e^{-0.0699x}$ & 42  & 0.9325 \\
\hline
\end{tabular}
\end{center}

Rewriting each in the $y = a \cdot R^x$ form:
\begin{align*}
\text{6 sided:} \quad & y = 136(0.8344)^x \\
\text{4 sided:} \quad & y = 44.1(0.7573)^x \\
\text{8 sided:} \quad & y = 32.8(0.8772)^x \\
\text{12 sided:} \quad & y = 33.4(0.9325)^x
\end{align*}

\begin{problem}
What do the $x$-value, the $y$-value, the $a$-value, and the $R$-value represent across
all of these models?

\begin{freeResponse}
\end{freeResponse}

\begin{solution}
The $x$-values represent the number of rolls.

The $y$-values represent the number of dice remaining after $x$ rolls.

The $a$-values represent the starting number of dice.

The $R$-values represent the percentage of dice remaining after a single roll, which
follows from the dice decreasing by a certain percentage after each roll.

It is important to specify that this is the percentage remaining after ONE roll. The
$R$-value is not the fraction remaining at the end, nor over the whole experiment ---
it is what gets applied once per roll.
\end{solution}
\end{problem}

\section*{Exploring Exponential Modelling}

Excel and Google Sheets give exponential models in the form $y = a \cdot e^{bx}$, and we
rewrote them as $y = a \cdot R^x$ to see the connection to the context. But here is the
striking part: \emph{without rolling any dice at all}, and without any special program, we
could have written a very decent model ourselves. Let's see how.

Do a thought experiment with 6-sided dice, following the same procedure: roll, and remove
any die that came up ``1.''

\begin{problem}
Theoretically, what fraction of the dice should roll a ``1'' for a 6-sided die? And what
fraction should \emph{not} roll a ``1''?

\begin{freeResponse}
\end{freeResponse}

\begin{solution}
A 1 in 6 chance, so $\tfrac{1}{6}$ likelihood of rolling a 1.

Therefore a 5 in 6 chance of not rolling a 1, since
$1 - \tfrac{1}{6} = \tfrac{5}{6}$.
\end{solution}
\end{problem}

\begin{problem}
Suppose we start with the class's actual starting number of 6-sided dice, 129. If we write
a theoretical model in the form $y = a \cdot R^x$, which value should 129 go in for, and
why? And which value should go in as $R$, and why?

\begin{freeResponse}
\end{freeResponse}

\begin{solution}
The 129 goes in for the $a$-value, since $a$ represents the starting number of dice and in
this thought experiment our starting number is 129.

The $\tfrac{5}{6}$ goes in for the $R$-value, since $R$ represents how many dice remain
after each roll --- which we get from how many are ``decaying'' away with each roll,
$1 - \tfrac{1}{6} = \tfrac{5}{6}$.

This is the key insight of the whole activity: the decay rate is not something you have to
measure. It is just the probability that a die survives a roll.
\end{solution}
\end{problem}

\begin{problem}
Put these numbers where they belong to create the theoretical model for our 6-sided dice.
Then compare its $R$ against the program's model, $y = 136(0.8344)^x$. Why would we
hope the program's $R$ is close to $0.8333\ldots$?

\begin{freeResponse}
\end{freeResponse}

\begin{solution}
The theoretical model is
\[
y = 129 \left(\frac{5}{6}\right)^x.
\]

The program's model has $R = 0.8344$, and $\tfrac{5}{6} = 0.8333\ldots$ --- a difference of
only $0.0011$.

We want the program's $R$ to land near $0.8333\ldots$ because that is what should
theoretically happen when these dice are rolled. Agreement between the fitted value and
the theoretical one indicates ``good data'' --- that the experiment behaved the way
probability says it should. A large gap would suggest either unusual luck or something
wrong with how the rolls were carried out or recorded.
\end{solution}
\end{problem}

\section*{Comparing the Two Models}

Continuing with the example data, which had 129 actual starting dice:

\begin{center}
\begin{tabular}{cccc}
\hline
Roll & Actual from experiment & Spreadsheet model $136(0.8344)^x$ & Theoretical model $129\left(\tfrac{5}{6}\right)^x$ \\
\hline
0 & 129 & 136 & 129 \\
1 & 111 & 113 & 108 \\
2 & 94  & 95  & 90  \\
3 & 84  & 79  & 75  \\
4 & 71  & 66  & 62  \\
\hline
\end{tabular}
\end{center}

\begin{problem}
How do these values compare? Which model is better, and which would we \emph{expect} to be
better? Why?

\begin{freeResponse}
\end{freeResponse}

\begin{solution}
The theoretical model is closer for the first two rolls --- at roll 0 it is exact, since it
was built from the true starting count. But from roll 2 onward the spreadsheet's model is
closer to the actual values, and would likely stay closer for many rolls after.

We would expect the spreadsheet model to be better for \emph{our} data, because it was
built from our data. Even though the theoretical model wins at rolls 0 and 1, the
spreadsheet finds the best exponential function for the data as a whole rather than for
any specific point, which is why it tracks better across the rest.

The theoretical model has a different kind of value. It would not change at all for a
different class rolling different 6-sided dice --- only the starting count $a$ would
differ. It describes what 6-sided dice \emph{do}, not what this particular pile happened
to do.
\end{solution}
\end{problem}

\section*{Theoretical Models for the Other Dice}

\begin{problem}
Using the same idea, write theoretical models for the 4, 8, and 12 sided dice, using each
group's actual starting count (44, 35, and 42 respectively). Also write one for 10-sided
dice, using $a$ as the starting number.

For $n$-sided dice, $R = \answer[tolerance=0.001]{0.9}$ when $n = 10$.

\begin{freeResponse}
\end{freeResponse}

\begin{solution}
For $n$-sided dice, exactly one face is a ``1,'' so the chance of surviving a roll is
$\tfrac{n-1}{n}$. That is the $R$-value every time.
\begin{align*}
\text{4-sided:} \quad & y = 44\left(\tfrac{3}{4}\right)^x \\
\text{8-sided:} \quad & y = 35\left(\tfrac{7}{8}\right)^x \\
\text{12-sided:} \quad & y = 42\left(\tfrac{11}{12}\right)^x \\
\text{10-sided:} \quad & y = a\left(\tfrac{9}{10}\right)^x
\end{align*}
Notice that the 10-sided model needs no experiment whatsoever --- we can write it down
without anyone ever having rolled a 10-sided die, because the rate follows from the
geometry of the die alone.

Note also the pattern: more sides means a larger $R$, closer to 1, meaning slower decay.
This makes sense --- with 12 sides, a die has only a $1/12$ chance of being removed each
roll, so the pile shrinks far more slowly than a 4-sided pile does.
\end{solution}
\end{problem}

\begin{problem}
Compare the theoretical models against the program's models:

\begin{center}
\begin{tabular}{lll}
\hline
Dice & Theoretical & Program's \\
\hline
6 sided  & $y = 129\left(\tfrac{5}{6}\right)^x$, \; $R = 0.8333$   & $y = 136(0.8344)^x$ \\
4 sided  & $y = 44\left(\tfrac{3}{4}\right)^x$, \; $R = 0.7500$    & $y = 44.1(0.7573)^x$ \\
8 sided  & $y = 35\left(\tfrac{7}{8}\right)^x$, \; $R = 0.8750$    & $y = 32.8(0.8772)^x$ \\
12 sided & $y = 42\left(\tfrac{11}{12}\right)^x$, \; $R = 0.9167$  & $y = 33.4(0.9325)^x$ \\
\hline
\end{tabular}
\end{center}

Which dice seemed to have ``acted'' the most theoretically, and which the least? Explain
how you decided.

\begin{freeResponse}
\end{freeResponse}

\begin{solution}
First lay out both differences for each:

\begin{center}
\begin{tabular}{lcc}
\hline
Dice & $|R$ difference$|$ & $|a$ difference$|$ \\
\hline
6 sided  & 0.0011 & 7.0 \\
4 sided  & 0.0073 & 0.1 \\
8 sided  & 0.0022 & 2.2 \\
12 sided & 0.0158 & 8.6 \\
\hline
\end{tabular}
\end{center}

\textbf{Least theoretical: the 12-sided dice}, and this one is unambiguous --- they are
furthest off on \emph{both} measures, with the largest $R$ gap (0.0158) and the largest
$a$ gap (8.6).

\textbf{Most theoretical: it depends which measure you use,} and this is worth being
careful about. On the $a$-value the 4-sided dice win easily, off by only 0.1 die. But on
the $R$-value the 6-sided dice win, off by only 0.0011 against the 4-sided dice's 0.0073
--- nearly seven times closer.

So a defensible answer names the criterion. If you care most about the decay rate matching
theory --- which is arguably the more meaningful comparison, since $R$ is the part that
probability predicts --- the 6-sided dice acted most theoretically. If you care about the
model reproducing the starting count, the 4-sided dice did.

Note that the $a$-value is not really a test of whether the dice behaved theoretically at
all. It measures how well the fitted curve happens to pass through roll 0, which is more a
feature of the curve-fitting than of the dice.
\end{solution}
\end{problem}

\end{document}
